\section{Cardinality}
\label{sec:cardinali}

\begin{definition}[Cardinality]
  We will say that two sets $A$ and $B$ have the same \emph{cardinality}%
\mymargin{cardinality}%
\index{cardinality} 
  (or are \emph{equipotent}) 
  \index{equipotence}%
  \index{sets!equipotent}%
  \index{set!equipotent}%
  if there exists a bijective function $f\colon A \to B$.
  In this case, we will write:
  \[
    \# A = \# B.  
  \] 
  If there exists an injective function $f\colon A\to B$, it means that 
  $A$ has the same cardinality as a subset of $B$ (since $f\colon A \to f(A)$
  is bijective). In this case, we will write:
  \[
    \# A \le \#B.
  \]
\end{definition}

Intuitively, if two sets have the same cardinality, 
it means they have the same number of elements.
However, the previous definition captures this concept without resorting 
to the notion of number. 
This has the advantage of making this definition applicable 
to any set, even those with \emph{infinite} elements.

Note that we have not provided a definition of $\#A$, and thus we are not 
defining what the cardinality of a set is, but we are merely defining 
a relation between sets that 
we denote, improperly, using an equality: $\#A = \#B$.

It is obvious that $\#A = \#A$ since the identity $\id_A$ is bijective.
It is also clear that if $\#A = \#B$ and $\#B = \#C$, then $\#A = \#C$ because 
composing two bijective functions results in another bijective function. 
Finally, if $\#A = \#B$, then $\#A \le \#B$ (since bijective functions 
are injective), and it is clear that if $\#A \le \#B$ and $\#B \le \#C$, then 
$\#A\le \#C$ (since the composition of injective functions is injective).

We can define the inverse symbol $\#A \ge \#B$ 
as $\#B\le \#A$.
Obviously, if $A\subset B$, then $\#A \le \#B$.
If $B\neq \emptyset$, 
the following theorem tells us that $\#A \ge \#B$ 
is equivalent to saying that 
there exists $f\colon A \to B$
surjective. 
%
\begin{theorem}\label{th:95444}
  Let $B\neq \emptyset$.
  There exists a surjective function $f\colon A\to B$ 
  if and only if there exists an injective function $g\colon B\to A$.
\end{theorem}
% 
\begin{proof}
On one hand, if there exists an injective function $g\colon B\to A$, then $g$ is
a bijection
between $B$ and $g(B)$. The inverse function $f\colon g(B) \to B$ 
\mynote{This function $f$ is called 
\emph{left inverse}
\index{inverse!left}%
of $g$ since $f(g(y))=y$ for every $y\in B$}
can be extended to all of $A$ by fixing an arbitrary value 
(this can be done if $B$ is not empty) at the points of $A\setminus g(B)$,
thus obtaining a surjective function from $A$ to $B$.
On the other hand, if there exists a surjective function $f\colon A\to B$,
for every $b\in B$, the set $f^{-1}(\ENCLOSE{b})$ is never 
empty, and thus it is clear that there must exist 
a function $g\colon B\to A$ such that $g(b)$ is any element 
of this set
\mynote{This function $g$ is called 
\emph{right inverse} 
\index{inverse!right}%
of $f$ 
since $f(g(y))=y$ for every $y\in B$}
(here the axiom~\ref{axiom:AC} discussed below is used). 
Clearly, $g$ is injective.
\end{proof}

In the proof of the previous theorem, we used the following axiom 
of set theory, which, surprisingly, is not a consequence of the axioms 
we have introduced so far.

\begin{axiom}[of choice]%
  \label{axiom:AC}%
  \index{axiom!of choice}%
  \index{choice!axiom of}%
  \index{AC}%
  Let $F\colon A \to \mathcal P(B)$
  be a function such that $F(a)\neq \emptyset$ 
  for every $a\in A$. Then there exists a function
  $f\colon A \to B$ such that $f(a)\in F(a)$
  for every $a\in A$.
\end{axiom}

The axiom of choice (often denoted by \emph{$AC$}%
\mymargin{$AC$}%
\index{AC}, \emph{axiom of choice})
is an axiom that is somewhat controversial,
and some mathematicians prefer not to use it in their proofs.
The formal system that defines set theory without 
introducing the axiom of choice 
is called \emph{$ZF$}%
\mymargin{$ZF$}%
\index{ZF} (Zermelo-Fraenkel), while 
if the axiom of choice is added, the theory is called 
\emph{$ZFC$}%
\mymargin{$ZFC$}%
\index{ZFC}.

Thanks to the axiom of choice, it is also possible to prove 
that given two sets $A$ and $B$, their cardinalities
can be compared: $\#A\le \#B$ or 
$\#B\le \#A$. 
However, the proof becomes complicated and goes beyond our purposes.

\begin{theorem}
  \label{th:cardinale_totale}%
  If $A$ and $B$ are any sets, then $\#A\le \#B$ or $\#B\le \#A$.
\end{theorem}

The following theorem instead demonstrates the antisymmetric property 
of the relation between cardinalities. 

\begin{theorem}
  If $\#A \le \#B$ and $\#B \le \#A$, then $\#A = \#B$.
\end{theorem}
\begin{proof}
This is none other than the Cantor-Bernstein theorem~\ref{th:cantor_bernstein}.
\end{proof}
%
Recall theorem~\ref{th:cardinali_finiti}
which tells us that $A$ is a finite set if and only if 
there exists $n\in\NN$ 
such that $\#A = \#[n]$, where we have set $[n]=\ENCLOSE{k\in \NN, k<n}
=\ENCLOSE{0,\dots,n-1}$.
Now we can observe that $\#A=n$ is an abbreviation of $\#A = \#[n]$.

The cardinality of $\NN$ is sometimes denoted by the symbol
\index{$\aleph_0$}%
\index{aleph}%
$\aleph_0$, thus we can write  
\[
    \# A = \aleph_0, \qquad
    \#A \le \aleph_0, \qquad
    \# A\ge \aleph_0
\]
to mean respectively 
\[
  \# A = \# \NN, \qquad
  \# A \le \# \NN, \qquad
  \# A \ge \# \NN.  
\]
\mynote{%
The symbol $\aleph$ is the first letter of the Hebrew alphabet 
and is pronounced \emph{aleph}. 
\\
Recall (theorem~\ref{th:unicitaN}) that the sets 
$\NN$ that satisfy the Peano axioms are all isomorphic to each other 
and therefore, in particular, all have the same cardinality.
We have already observed (theorem~\ref{th:esistenza_naturali}) 
that if $A$ is an infinite set, there exists $N\subset A$
that satisfies the Peano axioms. 
The inclusion $N\to A$ is injective, and thus $\aleph_0 = \# \NN \le \# A$
for any infinite $A$.
In this sense, $\NN$ is the \emph{smallest} infinite set, so $\aleph_0$ is the
smallest
infinite cardinality.
\\
It could also be shown that there must exist a cardinality $\aleph_1>\aleph_0$
which is the smallest cardinality greater than $\aleph_0$.
We know that $\# \mathcal P(\NN)>\# \NN$, but curiously, it is not possible to
prove
that $\aleph_1 = \# \mathcal P(\NN)$ (this is called the \emph{continuum
hypothesis}
\index{hypothesis!of the continuum}%
\index{continuum!hypothesis of the}%
since we will see in theorem~\ref{th:cantor_secondo} that $\# \mathcal
P(\NN)=\#\RR$
and the set $\RR$ is called \emph{continuum} whereas $\NN$ 
is called \emph{discrete}).
Thus, the continuum hypothesis is an additional axiom that could be 
added to the axioms of set theory: 
there are no sets with cardinality strictly between 
the cardinalities of $\NN$ and $\mathcal P(\NN)$.
}%

We have already observed that $\NN$ is an infinite set, and thus 
also $\ZZ$, $\QQ$, and $\RR$ are infinite.
For finite sets, if $A\subset B$ but $A\neq B$, it results in $\#A < \# B$
since $\#A - \#B = \#(B\setminus A) > 0$. 
However, we must observe that the same does not hold for infinite sets.
In fact, the function $s\colon \NN \to \NN$ defined by $s(x)=x+1$
is a bijection between $\NN$ and $\NN\setminus\ENCLOSE{0}$. 
Thus $\#\NN = \#(\NN\setminus\ENCLOSE{0})$ 
despite the difference between the two sets $\ENCLOSE{0}$ having 
cardinality $1$. 
Even the set of even numbers or the set 
of perfect squares (as Galileo had already noted)
have the same cardinality as $\NN$ since the functions $n\mapsto 2n$
and the function $n\mapsto n^2$ are injective.

It is not difficult to find an injective function from $\ZZ$ to $\NN$
(we leave this as an exercise): this demonstrates that $\#\ZZ = \# \NN$ as well.

Cantor even demonstrates that $\#\QQ = \#\NN$: 
to do this, he uses the following.

\begin{theorem}[Cantor's First Diagonal Method]
\label{th:Cantor_primo}%
\index{Cantor!first diagonal method}%
The set $\NN \times \NN$ has the same cardinality as $\NN$. Consequently,
\[
  \# \NN = \# \ZZ = \# \QQ.
  \]
\end{theorem}
%
\begin{proof}[Idea of the proof]
  Number the squares of an infinite chessboard
  as in Figure~\ref{fig:cantor1}.
  
  Since it is quite easy to find an injective function 
  from $\QQ$ to $\NN \times \ZZ$, it is easily deduced that $\#\QQ = \#\NN$.
\end{proof}

\begin{figure}
  \begin{tabular}{c|ccccccc}
   $m$ & $\ddots$ & $\ddots$ & $\ddots$ & $\ddots$ & $\ddots$ & $\ddots$\\
   $\vdots$ & $\ddots$ & $\ddots$ & $\ddots$ & $\ddots$ & $\ddots$ & $\ddots$\\
   3 & 9 & $\nwarrow$ & $\ddots$  & $\ddots$ & $\ddots$ & $\ddots$\\
   2 & 5 & 8 & 12 & $\ddots$  & $\ddots$ & $\ddots$\\
   1 & 2 & 4 & 7 & 11 & $\ddots$  & $\ddots$\\
   0 & 0 & 1 & 3 & 6 & 10 & $\ddots$ \\ \hline
     & 0 & 1 & 2 & 3 & 4 & $\dots$ & $n$
  \end{tabular}
  \caption{
    The diagonal numbering of the squares
    of an infinite chessboard. It could be verified
    that the number present in the square with coordinates $(n,m)$
    is written as $f(n,m) = \frac{(n+m+1)(n+m)}{2}+m$
    and is a bijective function $f\colon \NN\times\NN \to\NN$.}
  \label{fig:cantor1}
\end{figure}

At this point, one might think that
all infinite sets are countable.
However, given any set (even infinite),
there always exists a set with greater cardinality,
as demonstrated in the following theorem, which is a small gem of logic.
In fact, the method used is akin to the liar paradox 
and is the same idea used in Russell's paradox (theorem~\ref{th:Russell}).
\mynote{The set $\mathcal P(A)$ is the power set, 
defined on page~\pageref{def:insieme_parti}}%
%
\begin{theorem}[Cantor]%
\index{Cantor!cardinality of the power set}%
\index{Theorem of!Cantor}%
\label{th:Cantor}%
  If $A$ is any set, then $\# \mathcal P(A) > \# A$.
\end{theorem}
%
\begin{proof}
  It is clear that $\# A \le \#\mathcal P(A)$ since 
  the function $f\colon A \to \mathcal P(A)$ defined by $f(x)=\ENCLOSE{x}$
  is obviously injective.

    Suppose, for contradiction, that there exists a bijective function $f\colon
  A\to \mathcal P(A)$
  and consider the set 
  \[
    C = \ENCLOSE{x\in A \colon x\not\in f(x)}.  
  \]
  If $f$ were surjective, there should exist $c\in A$ such that $f(c) = C$.
  We can then ask whether $c\in C$ and discover that, 
  by definition of $C$, the proposition $c\in C$ is equivalent 
  to $c\not\in f(c) = C$. 
  Thus $c\in C \iff c\not\in C$, which is absurd.
\end{proof}
%
\begin{corollary}[there does not exist the set of all sets]
  If there existed a set $\mathcal U$ (universe) 
  such that $\forall x\colon x \in U$, then 
  we would have $\mathcal \P(\mathcal U) \subset \mathcal U$
  contradicting the previous theorem.
\end{corollary}
%
Regarding the set $\RR$, it
turns out that $\#\RR > \#\NN$
(again thanks to Cantor, theorem~\ref{th:cantor_secondo} below)
and it could be demonstrated that indeed $\#\RR = \#\mathcal P(\NN)$.

\begin{theorem}[Cantor's Second Diagonal Method]
  \label{th:cantor_secondo}
The Cantor set
\[
  C = \{r(\vec a)\colon \vec a \in \ENCLOSE{0,2}^\NN\}
\]
has cardinality
\[
\# C = \# \mathcal P(\NN).
\]
In particular, $\#\RR > \#\NN$.
\end{theorem}
%
\begin{proof}
Consider the function $r$ defined in chapter~\ref{sec:decimali} that associates 
to a sequence of binary digits the corresponding real number.
We want to verify that $r\colon \ENCLOSE{0,2}^\NN \to \closeinterval{0}{1}$
is injective.
We have already seen in general that $r(\vec a) = r(\vec b)$ only 
if the first differing digit in $\vec a$ and $\vec b$ differs by one unit.
But since all digits of $\vec a$ and $\vec b$ are by hypothesis $0$ or 
$2$, they do not differ by one unit, and thus $r$ is injective.
Therefore, $\# C = \#\ENCLOSE{0,2}^\NN$.
On the other hand, $\#\ENCLOSE{0,2}^\NN = \#\mathcal P(\NN)$ since 
each $\vec a \in \ENCLOSE{0,2}^\NN$ can be put in a one-to-one correspondence 
with the set $\vec a^{-1}(\ENCLOSE{0})$ of indices $k\in \NN$
for which $a_k=0$.

From theorem~\ref{th:Cantor}, we deduce that $\# C = \#\mathcal P(\NN) > \#\NN$ 
and since $C\subset \RR$, it follows that $\# \RR \ge \# C >\#\NN$.

On the other hand, we can show that $\#\RR \le \# \mathcal P(\NN)$
since we can construct a function $f\colon \RR \to \mathcal P(\QQ)$
in this way:
\[
f(x) = \ENCLOSE{q\in \QQ\colon q<x}.
\]
Due to the density of $\QQ$ in $\RR$, we know that if $x<y$, 
there exists $q\in \QQ$ with $x<q<y$, and thus $q\in f(y)\setminus f(x)$ 
from which $f(x)\neq f(y)$. 
Thus $f$ is injective, which means $\#\RR\le \# \mathcal P(\QQ)$.
But $\#\NN = \#\QQ$, so $\#\mathcal P(\QQ) = \#\mathcal P(\NN)$, 
and the proof is concluded.
\end{proof}

The set $C$ defined in the previous proof with $d=3$ 
is the same \emph{Cantor set}%
\mymargin{Cantor set}%
\index{set!of Cantor}
\index{Cantor!set of}% 
considered in example~\ref{ex:insieme_Cantor}.
In fact, it could be easily shown that 
$C=\frac 1 3 C \cup (\frac 2 3 + \frac 1 3 C)$.